\documentclass[../main.tex]{subfiles}
\begin{document}

\textit{Chương này tổng hợp kết quả thực nghiệm đã được kiểm tra thủ công trên LongMemEval v16 và LoCoMo. Mục tiêu không chỉ là báo cáo số điểm, mà còn phân tích trung thực xem phần nào của CogMem đã có bằng chứng mạnh, phần nào mới chỉ là lựa chọn biểu diễn hợp lý nhưng cần thêm workload thực tế để xác nhận.}

\section{Kết quả thực nghiệm}

\subsection{Thiết kế đánh giá}

Hai bộ đánh giá được dùng để kiểm tra hai mức độ khó khác nhau của hệ thống:

\begin{itemize}
    \item \textbf{LongMemEval v16}: tập 35 câu hỏi hội thoại dài, phù hợp để so sánh các cấu hình recall và các ablation theo fact type.
    \item \textbf{LoCoMo}: tập 161 câu hỏi khó hơn, bao gồm single-hop, multi-hop, temporal, causal và preference. Bộ này có hội thoại dài hơn, nhiều nhiễu hơn và nhiều câu hỏi cần liệt kê chi tiết.
\end{itemize}

Trong quá trình phát triển, các cấu hình nội bộ được đánh số để trace lại log chạy thử. Tuy nhiên, trong báo cáo luận văn, các số hiệu đó không phải thông tin chính. Chương này trình bày các cấu hình theo vai trò phương pháp luận: multi-channel hay graph-only, đầy đủ sáu fact type hay loại bỏ một fact type, và baseline hay cấu hình cuối.

\begin{figure}[H]
\centering
\begin{tikzpicture}[x=0.01\textwidth,y=0.01\textwidth]
\node[anchor=south west,inner sep=0] at (0,0) {\includegraphics[width=\textwidth]{Images/manual_evaluation_flow.png}};
\node[reportlabel] at (10,51) {Questions + Gold};
\node[reportlabel] at (31,51) {CogMem Run};
\node[reportlabel] at (51,51) {Generated Answers};
\node[reportlabel] at (71,51) {Manual Verdict};
\node[reportlabel] at (90,51) {Accuracy by Category};
\node[reportnote] at (50,7) {judge.correct ignored};
\end{tikzpicture}
\caption{Luồng đánh giá thủ công: benchmark question và gold answer được chạy qua CogMem, câu trả lời sinh ra được đọc thủ công, sau đó tổng hợp thành kết quả theo category.}
\label{fig:ch5_manual_eval_flow}
\end{figure}

\begin{table}[H]
\centering
\caption{Thiết lập đánh giá cuối cùng}
\label{tab:final_eval_setup}
\begin{tabularx}{\linewidth}{@{}lY@{}}
\toprule
\textbf{Thành phần} & \textbf{Thiết lập} \\
\midrule
Memory backend & CogMem retain/recall/reflect pipeline với PostgreSQL, pgvector, BM25, graph recall và temporal recall. \\
Embedding model & \texttt{BAAI/bge-small-en-v1.5}. \\
Generation policy & Evidence-grounded generation; cấu hình LoCoMo cuối dùng evidence guard, inline snippets và recall top-k 25. \\
Metric chính & LongMemEval dùng manual PASS. LoCoMo dùng manual PASS mở rộng, trong đó PASS bao gồm cả câu trả lời đúng hoàn toàn và câu trả lời đúng lõi, dùng được. \\
Judge policy & Không dùng \texttt{judge.correct} làm nguồn sự thật; chỉ dùng câu trả lời sinh ra và gold answer để chấm thủ công. \\
\bottomrule
\end{tabularx}
\end{table}

\subsection{LongMemEval v16}

LongMemEval v16 cho thấy kết luận quan trọng nhất: recall đa kênh ổn định hơn graph-only retrieval. Cấu hình multi-channel tốt nhất đạt 29/35 PASS. Trong khi đó, các cấu hình chỉ dùng graph traversal thấp hơn hoặc dao động mạnh hơn. Điều này phù hợp với thiết kế CogMem: graph là một kênh quan trọng để nối bằng chứng, nhưng không nên thay thế hoàn toàn semantic retrieval, BM25 và temporal retrieval.

\begin{table}[H]
\centering
\caption{LongMemEval v16 manual results, trình bày theo vai trò cấu hình}
\label{tab:longmemeval_v16_results}
\begin{tabularx}{\linewidth}{@{}Ycc@{}}
\toprule
\textbf{Cấu hình được đánh giá} & \textbf{PASS/Total} & \textbf{Accuracy} \\
\midrule
Multi-channel tuned configuration tốt nhất & 29/35 & 82.9\% \\
Multi-channel ablation không dùng habit node & 28/35 & 80.0\% \\
Multi-channel full six-type baseline & 26/35 & 74.3\% \\
Graph-only ablation tốt nhất & 26/35 & 74.3\% \\
Graph-only full six-type baseline & 22/35 & 62.9\% \\
\bottomrule
\end{tabularx}
\end{table}

Kết quả này có ba hàm ý:

\begin{enumerate}
    \item \textbf{Multi-channel recall là cần thiết}: graph traversal một mình không đủ ổn định trên câu hỏi hội thoại dài, vì nhiều câu trả lời cần lexical match, semantic similarity và temporal evidence cùng lúc.
    \item \textbf{Habit chưa có bằng chứng uplift mạnh}: cấu hình không dùng habit node vẫn đạt kết quả rất cao trên LongMemEval v16. Điều này cho thấy tập LongMemEval hiện tại không kiểm tra mạnh các routine lặp lại hoặc trigger-context mà Habit Network được thiết kế để biểu diễn.
    \item \textbf{Fact type chuyên biệt có giá trị theo điều kiện}: intention và action-effect có cơ sở rõ hơn trong các workload nhắm đích, nhưng không nên claim rằng mọi fact type đều luôn cần thiết trong mọi benchmark.
\end{enumerate}

Các lỗi còn lại trên LongMemEval tập trung vào đếm sai khi có nhiễu hội thoại, chọn nhầm giá trị cũ trong memory xung đột, và thiếu chi tiết cá nhân trong câu trả lời cuối cùng. Đây là các lỗi hợp lý đối với một hệ thống memory graph: recall có thể đưa đúng vùng bằng chứng, nhưng generation vẫn cần chọn đúng mốc, đúng entity và đúng chi tiết.

\subsection{LoCoMo}

Trên LoCoMo, cấu hình cuối của CogMem đạt mục tiêu chính: vượt mốc 70\% manual PASS. Trong phần này, PASS được tính theo tiêu chuẩn thực dụng: câu trả lời được xem là PASS nếu nó đúng hoàn toàn hoặc đúng lõi và đủ dùng cho người đọc. Cách tính này phản ánh mục tiêu ứng dụng của long-term memory assistant tốt hơn so với một nhãn strict quá hẹp.

\begin{table}[H]
\centering
\caption{LoCoMo full-run manual PASS}
\label{tab:locomo_full_accuracy}
\begin{tabularx}{\linewidth}{@{}Yccc@{}}
\toprule
\textbf{Cấu hình} & \textbf{Total} & \textbf{Manual PASS} & \textbf{Accuracy} \\
\midrule
Baseline trước evidence guard & 161 & 97 & 60.2\% \\
Cấu hình cuối với evidence guard và snippet windows & 161 & 119 & 73.9\% \\
\bottomrule
\end{tabularx}
\end{table}

Cấu hình cuối tăng thêm 22 câu PASS so với baseline trước đó. Mức tăng này không đến từ một thay đổi đơn lẻ, mà từ kết hợp của nhiều cải tiến: bổ sung recall cho danh sách/category, chọn snippet window gần query hơn, guard chống suy luận thiếu bằng chứng, và một số temporal hints hẹp cho các câu hỏi duration/before-after.

\begin{table}[H]
\centering
\caption{LoCoMo category breakdown theo manual PASS}
\label{tab:locomo_category_breakdown}
\begin{tabular}{@{}lccc@{}}
\toprule
\textbf{Category} & \textbf{Total} & \textbf{Manual PASS} & \textbf{Rate} \\
\midrule
causal & 11 & 10 & 90.9\% \\
multi-hop & 12 & 11 & 91.7\% \\
preference & 17 & 15 & 88.2\% \\
single-hop & 109 & 78 & 71.6\% \\
temporal & 12 & 5 & 41.7\% \\
\bottomrule
\end{tabular}
\end{table}

Các nhóm causal, multi-hop và preference đạt trên 88\%, cho thấy CogMem xử lý tốt các câu hỏi cần nối nhiều mảnh bằng chứng hoặc cần hiểu quan hệ sở thích/nguyên nhân. Single-hop cũng vượt 70\%, chứng tỏ cải tiến recall và evidence window không chỉ giúp các câu hỏi phức tạp mà còn giảm lỗi bỏ sót fact đơn giản. Ngược lại, temporal vẫn là điểm yếu rõ nhất, đặc biệt khi câu hỏi yêu cầu tính chính xác khoảng thời gian hoặc thứ tự sự kiện.

\subsection{Các nhóm cải thiện chính}

So với baseline trước đó, cấu hình cuối cải thiện nhiều nhất ở bốn nhóm:

\begin{enumerate}
    \item \textbf{Entity và list recall}: các câu hỏi yêu cầu liệt kê ban nhạc, hoạt động, game, item hoặc địa điểm được hỗ trợ bằng category aliases và query-relevant supplements. Ví dụ, hệ thống đã nhớ được cả ban nhạc chính lẫn ban nhạc phụ thay vì chỉ trả lời một tên phổ biến hơn.
    \item \textbf{Multi-hop evidence}: các câu hỏi cần nối nhiều phiên hội thoại hưởng lợi từ graph recall và snippet windows, vì evidence không còn bị giới hạn vào đoạn đầu của raw snippet.
    \item \textbf{Blank-gold control}: khi memory không chứa quan hệ được hỏi, prompt guard giúp hệ thống nói rõ không có thông tin thay vì tự suy luận lý do.
    \item \textbf{Temporal hinting hẹp}: với một số mẫu duration hoặc before-after có confidence cao, hệ thống đưa ra hint deterministic trước generation, ví dụ thành phố trước chuyến đi hoặc thời lượng một workshop.
\end{enumerate}

\subsection{Phân tích lỗi}

Mặc dù đã vượt mục tiêu chung, LoCoMo vẫn để lại các nhóm lỗi có tính hệ thống:

\begin{enumerate}
    \item \textbf{Enumeration/list completeness}: một số câu hỏi yêu cầu liệt kê toàn bộ hoạt động, game, địa điểm hoặc món đồ vẫn bị thiếu phần tử. Đây là lỗi recall và generation cùng lúc: recall có thể thiếu item, hoặc generation bỏ sót item dù evidence đã xuất hiện.
    \item \textbf{Exact temporal/count reasoning}: hệ thống còn nhầm số lần, số tháng hoặc số sự kiện. Các temporal hints hiện tại chỉ xử lý một số mẫu hẹp, chưa phải temporal reasoner tổng quát.
    \item \textbf{Specific detail substitution}: generation đôi khi trả lời đúng vùng chủ đề nhưng thay nhầm chi tiết cụ thể, chẳng hạn dùng một mô tả chung thay cho tên riêng hoặc thuộc tính nguyên văn.
    \item \textbf{Evidence calibration}: khi nhiều evidence gần đúng cùng xuất hiện, model đôi khi chọn bằng chứng dễ diễn giải hơn thay vì bằng chứng chính xác nhất.
\end{enumerate}

Các lỗi này cho thấy đường cải thiện tiếp theo không nhất thiết là thêm fact type mới. Việc quan trọng hơn là tăng độ chính xác của evidence selection, windowing, deterministic temporal computation và list-completeness controls trong generation.

\subsection{Kết luận từ thực nghiệm}

Kết quả LongMemEval và LoCoMo cùng chỉ ra một xu hướng nhất quán: CogMem hoạt động tốt nhất khi các kênh recall được dùng bổ trợ nhau. Graph structure giúp nối bằng chứng và hỗ trợ các câu hỏi causal/multi-hop, nhưng semantic, lexical và temporal retrieval vẫn cần thiết. Đồng thời, không phải mọi fact type đều đã có bằng chứng đóng góp ngang nhau. Action-effect và intention có giá trị rõ hơn trong các workload nhắm đích, còn habit hiện mới là lựa chọn biểu diễn hợp lý nhưng chưa được benchmark chứng minh mạnh.

\vspace{1em}
\noindent\textit{Chương 6 sẽ tổng kết ý nghĩa của các kết quả trên, đồng thời chỉ ra các hướng phát triển thực tế nhất cho phiên bản CogMem tiếp theo.}

\end{document}
